%% =====================================================================
%%  7 Discussion
%% =====================================================================
%%  GUIDE (this section is where gold papers separate from silver ones)
%%   - Validation: compare with an external number (published experiment,
%%     flight data, a limiting case). Quantify the deviation; "good
%%     agreement" is not a number.
%%   - Error budget: a table with one row per source and a dominant term.
%%   - Sensitivity: what changes the conclusion, and by how much.
%%   - Limits: the regime where the model must not be used.
%%  Only 18 of 74 gold papers contain an explicit validation -- doing it
%%  well is a cheap differentiator.
%% =====================================================================
\section{Discussion}
\label{sec:discussion}

\subsection{Validation against published data}
\label{sec:disc-validation}

%% EXAMPLE: the comparison below is a template. Fit C_d to the data you
%% actually have, and cite the source you verified yourself.
The model has one free parameter, $C_d$. Fitting it to a published
leak-down curve for a comparable module gives $C_d=0.82$; with that value
the model reproduces the curve within \SI{8}{\percent} over
$1>p/p_0>0.35$. The largest deviation sits at the low-pressure end, where
the measured decay is slower than the adiabatic solution predicts, and it
is reduced from \SI{8}{\percent} to \SI{4}{\percent} once the wall term
\eqref{eq:temp} is switched on. That is the evidence that the thermal
correction is capturing something real rather than fitting noise.

A second, independent check is the limiting case of no wall heating. With
$\tau_w\to\infty$ the same code gives $5$ and \SI{20}{cm} times in the
ratio $3.99:1$, an exponent of \num{2.00}, which is the value
Eq.~\eqref{eq:tau} demands for $A\propto d^2$. The \num{1.97} of
Section~\ref{sec:results} is therefore produced by the thermal correction
and not by the integrator, and the two limits bracket the physical answer.

\subsection{Error budget}
\label{sec:disc-budget}

\begin{table}[htbp]
  \centering
  \small
  \caption{Error budget for $t_{0.3}$ at $d=\SI{10}{cm}$. The dominant
  term is the hole diameter, which enters as $d^{-2}$; the discharge
  coefficient is second.}
  \label{tab:budget}
  \begin{tabular}{@{}lS[table-format=2.0]S[table-format=2.1]l@{}}
    \toprule
    Source & {Magnitude (\%)} & {Effect on $t_{0.3}$ (\%)} & Handling \\
    \midrule
    Hole diameter         &  5 & 10.5 & geometric measurement \\
    Discharge coefficient &  6 &  6.3 & calibrated, Sec.~\ref{sec:disc-validation} \\
    Module volume         &  2 &  1.9 & module geometry \\
    Wall temperature      &  3 &  0.5 & lumped model, Eq.~\eqref{eq:temp} \\
    Numeric               & {--} &  0.09 & step halving \\
    \midrule
    Total (quadrature)    & {--} & 12.4 & dominant: $d$ \\
    \bottomrule
  \end{tabular}
\end{table}

Table~\ref{tab:budget} gives $\pm\SI{12}{\percent}$ in total, so the answer
to Q1 is quoted as \SI{286}{s} rather than \SI{285.5}{s}. The error is
dominated by a length that can be measured directly, not by the model, and
that is the distinction the section exists to make: the physics is worth
\SI{6}{\percent}, the measurement is worth \SI{10}{\percent}.

\subsection{Limits of validity}
\label{sec:disc-limits}

Two limits bound the result. The quasi-steady assumption fails once the
venting time approaches the acoustic transit time of the module,
$L/c_s\approx\SI{30}{ms}$; that is four orders of magnitude below
$t_{0.3}$ at \SI{10}{cm} but only a factor of twenty at \SI{1}{m}, so the
model must not be pushed to metre-scale breaches without a compressible
treatment. The well-mixed assumption fails for a compartmented module well
before the pressure does: once the breach is sealed off from part of the
volume, the problem becomes a network of reservoirs and this single-volume
model overestimates the available time.

There is also a modelling limit worth stating plainly. Because $C_d$
absorbs the throat density ratio and the contraction, its fitted value
($0.82$) is not comparable with the handbook value for a sharp-edged
orifice ($0.61$): the two differ by the factor
$(2/(\gamma+1))^{1/(\gamma-1)}=0.63$ that the lumped form hides. A reader
who substitutes the handbook number into \eqref{eq:mass} will get a time
that is \SI{34}{\percent} too long.
